% Main Preamble

\documentclass[12pt, letterpaper]{article}
\title{Advion ESI-MS Maintenance Guide\\[0.12in] \large First Edition}
\author{Caleb Griffin}
\date{October, 2026}

\usepackage{
	xcolor, 
	graphicx, 
	titlesec, 
	geometry, 
	listings, 
	amsmath,
	hyperref
}
\usepackage{fontawesome5}
\usepackage[default]{lato}
\usepackage[T1]{fontenc}
\geometry{left=1.0in, right=1.0in, top=1.0in, bottom=1.0in}

\pagecolor{white}

\hypersetup{
	colorlinks=true,
	linkcolor=black,
	filecolor=magenta,
	urlcolor=blue
}

\begin{document}
\maketitle
\thispagestyle{empty}

\pagebreak
\tableofcontents
\pagebreak

\section{Introduction}
This is a detailed guide for the maintenance and troubloshooting for the Advion Expression L ESI-MS.
This is not an official guide from Advion, but the ``tips and tricks'' from a graduate student in Dr. Alex Lippert's lab.

\section{General Knowledge on ESI-MS}
There are a few concepts that are best to understand before using the ESI-Ms proficiently. 

\subsection{Scan speed}

\subsection{Scan range}

\subsection{Quadrapole MS}

\subsection{Ionization}
Ionization charges, coefficient, hardness.

\subsection{}

\subsection{}

\subsection{}

\subsection{}

\subsection{}

\subsection{}


\section{Common Usage Guide}
Common usage of the ESI-MS is in seven stages:
\begin{itemize}
	\item Sample preperation.
	\item Instrument setup.
	\item Pre-run cleaning.
	\item sample injection.
	\item post-run cleaning.
	\item Instrument satndby.
	\item Data analysis.
\end{itemize}
Each of these steps needs to be performed for the instrument to read well and perform well over a long period of time.


\subsection{Sample preparation}

\subsection{Instrument setup}

\subsection{Pre-run cleaning}

\subsection{sample injection}

\subsection{post-run cleaning}

\subsection{Instrument satndby}

\subsection{Data analysis}

\section{Calibration}
Calibrating the instrument is performed at a regular interval, typically  or 14 days.
Calibration is performed independently for the positive and negative modes, but the same procedure is performed for both.
Calibration is done in the tune/calibration tab on the right side of the advion mass express app.

\subsection{Peak Standards}
These are the masses of the corresponding signals from the tuning solution.
There are different masses for the negative and positive modes and the tuning solution produces different species for both. You can remove and add peak masses.

\subsection{Running a calibration}

\subsection{Testing a calibration}

\subsection{Reading the calibration report}

\subsection{}

\section{Cleaning \& Maintenance}
Cleaning the instrument and the occasional upkeep tasks are done as needed and usually not on a schedule.

\subsection{Cleaning the ionization cone}

\subsection{Replenishing solvent}

\subsection{}

\subsection{}

\subsection{}


\section{What not to do}
There are multiple small things that can either damage the instrument or suppress signal.
This is a list of things to avoid when using or maintaining the ESI-MS.
\begin{itemize}
	\item \textbf{Bad solvents}\\
	Do not use high boiling point solvents like DMSO, DMF, or NMP for sample preparation or as the mobile soution.
	
	\item \textbf{No nitrogen flow}\\
	
	\item \textbf{Pulling on isocratic pump}\\
	
	\item \textbf{Shaking the instrument}\\
	
	\item \textbf{Using non-HPLC grade solvent}\\
	
	\item \textbf{}\\

\end{itemize}


\section{FAQ}
\begin{itemize}
	\item \textbf{Why is the tune not working for a single mass?}\\
		Sometimes a fragment mass for the runing solution is not bright on the tuning run.
		Remove the mass from the calibration proceedure for preliminary calibration and return the mass at a later calibration.
	\item \textbf{How often should I tune the instrument?}\\
		Once every 7 or 14 days is typically good practice.
		
	\item \textbf{Why are my sample peaks negative on the counts/time graph?}\\
		Not sure, but I have occasionally resolved this issue by simply diluting the sample more.
		It does appear that this sometimes is just a characteristic of a poor analyte that may not be studied well by ESI-MS.
\end{itemize}

\end{document}
